% The dependencies that are real, and nothing else.
\begin{tikzpicture}[font=\sffamily\scriptsize, x=1cm, y=1cm,
  ch/.style={blk, text width=30mm, minimum height=8mm, inner sep=2pt,
             font=\sffamily\tiny},
  dep/.style={-Latex, draw=linecol, line width=0.6pt}]

% ================================================================ part I
\node[chlabel, anchor=west, fill=white] at (-7.4,3.9) {the node comes up};
\node[ch, fill=hwcol] (c1) at (-6.0,3.2) {1 what a joint node is};
\node[ch, fill=hwcol] (c2) at (-2.4,3.2) {2 the control period};
\node[ch, fill=hwcol] (c3) at (1.2,3.2) {3 one clock};
\node[ch, fill=hwcol] (c4) at (4.8,3.2) {4 the board support package};
\draw[dep] (c1) -- (c2);
\draw[dep] (c2) -- (c3);
\draw[dep] (c3) -- (c4);

% ================================================================ part II
\node[chlabel, anchor=west, fill=white] at (-7.4,2.3) {sensing and acting};
\node[ch, fill=usrcol] (c5) at (-6.0,1.6) {5 the encoder};
\node[ch, fill=usrcol] (c6) at (-2.4,1.6) {6 the inertial unit};
\node[ch, fill=usrcol] (c7) at (1.2,1.6) {7 the actuator you do not have};
\node[ch, fill=usrcol] (c8) at (4.8,1.6) {8 force and torque};
\draw[dep] (c4.south) -- (c8.north);

% ================================================================ part III
\node[chlabel, anchor=west, fill=white] at (-7.4,0.7) {the bus};
\node[ch, fill=kercol] (c9)  at (-6.0,0.0) {9 CAN-FD from the controller out};
\node[ch, fill=kercol] (c10) at (-2.4,0.0) {10 the motion master};
\node[ch, fill=kercol] (c11) at (1.2,0.0) {11 a joint protocol};
\node[ch, fill=kercol] (c12) at (4.8,0.0) {12 network management};
\node[ch, fill=kercol] (c13) at (4.8,-1.2) {13 two speeds on one wire};
\draw[dep] (c5.south) -- (c9.north);
\draw[dep] (c9) -- (c10);
\draw[dep] (c10) -- (c11);
\draw[dep] (c11) -- (c12);
\draw[dep] (c12) -- (c13);

% ================================================================ part IV
\node[ch, fill=fwcol] (c14) at (-6.0,-2.6) {14 a middleware participant};
\node[ch, fill=fwcol] (c15) at (-2.4,-2.6) {15 state and command as messages};
\node[ch, fill=fwcol] (c16) at (1.2,-2.6) {16 the loop closed over the bus};
\node[ch, fill=extcol] (c17) at (4.8,-2.6) {17 a real-time fieldbus, refused};
\draw[dep] (c9.south) -- (c14.north);
\draw[dep] (c13.south) -- (c17.north);
\draw[dep] (c14) -- (c15);
\draw[dep] (c15) -- (c16);
% drawn after its arrows so the label sits on top of them rather than under
\node[chlabel, anchor=west, fill=white] at (-7.4,-1.9) {the system above};

% the two that make chapter 16 possible
\draw[dep, draw=blue!55!black] (c7.east) -- (3.0,1.6) -- (3.0,-2.6) -- (c16.east);
\draw[dep, draw=blue!55!black] (c5.west) -- (-7.8,1.6) -- (-7.8,-3.6) -- (1.2,-3.6) -- (c16.south);
\node[chlabel, anchor=west, text=blue!55!black] at (-7.3,-3.95)
  {chapter 5's encoder and chapter 7's plant model are what make chapter 16's loop close at all};

% ================================================================ part V
\node[chlabel, anchor=west, fill=white] at (-7.4,-4.5) {making it trustworthy};
\node[ch, fill=red!18] (c18) at (-6.0,-5.2) {18 safe states};
\node[ch, fill=red!18] (c19) at (-2.4,-5.2) {19 update over the bus};
\node[ch, fill=red!18] (c20) at (1.2,-5.2) {20 the rig};
\draw[dep] (c18) -- (c19);
\draw[dep] (c19) -- (c20);
\draw[dep] (c16.south) -- (1.2,-4.2) -- (2.4,-4.2) -- (2.4,-5.2) -- (c20.east);

\node[umlnote, text width=52mm, anchor=north west] at (-7.4,-6.1)
  {Chapter 1 gates everything. Chapter 2's period and chapter 3's clock are
   assumed by every chapter after them and are not drawn again. Chapter 9's bus
   gates chapters 10 to 14 and 19.};
\node[umlnote, text width=52mm, anchor=north west] at (-1.4,-6.1)
  {Chapters 17 and 20 depend on almost everything and are the two that step back
   and look at the whole. Everything not drawn here is a preference rather than a
   requirement, which is most of the order.};
\end{tikzpicture}
